\section{Planteamiento del Problema, Hipótesis y Objetivos}
\label{sec:planteamiento}

\subsection{Planteamiento del Problema}
A pesar del enorme potencial de las heteroestructuras formadas por aleaciones de Heusler y materiales bidimensionales, la comprensión a nivel atómico y electrónico de sus interfaces sigue siendo limitada. En particular, la transferencia de carga en la interfaz, el desajuste de red (lattice mismatch), las reconstrucciones estructurales y la presencia de grupos funcionales en el MXene pueden alterar significativamente la estructura de bandas electrónicas, destruyendo o alterando los estados topológicos deseados. Por lo tanto, el problema científico radica en identificar cómo las características estructurales y químicas de la interfaz Heusler/MXene dictan las propiedades topológicas del sistema compuesto y determinar los mecanismos eficientes para su modulación y control externo.

\subsection{Pregunta de Investigación}
¿De qué manera influyen el ordenamiento atómico de la interfaz, el acoplamiento espín-órbita y los efectos de tensión epitaxial en la emergencia y modulación de fases topológicas no triviales en heteroestructuras compuestas de Heusler/MXene?

\subsection{Justificación}
El desarrollo de la espintrónica y la computación cuántica topológica requiere materiales con estados electrónicos altamente estables y protegidos contra la retrodispersión. Las heteroestructuras Heusler/MXene representan una nueva frontera para el diseño de materiales a medida. Comprender los principios físicos fundamentales que rigen estas interfaces mediante cálculos de primeros principios permitirá acelerar el diseño de dispositivos electrónicos sin disipación de calor, reduciendo el costo experimental y guiando la síntesis de nuevos materiales funcionales.

\subsection{Hipótesis}
La ingeniería de interfaces en heteroestructuras formadas por aleaciones de Heusler y MXenes modificados superficialmente permite inducir y sintonizar estados topológicos de borde no triviales mediante la transferencia selectiva de carga y el acoplamiento magnético interfacial, logrando una modulación dinámica a través de la aplicación de deformación uniaxial/biaxial y campos eléctricos externos.

\subsection{Objetivos}

\subsubsection{Objetivo General}
Estudiar teóricamente la estabilidad estructural, las propiedades electrónicas y la topología de bandas de las heteroestructuras Heusler/MXene mediante simulaciones de primeros principios basadas en DFT, con el fin de proponer mecanismos eficientes para la modulación y el control de fases topológicas.

\subsubsection{Objetivos Específicos}
\begin{enumerate}
    \item Diseñar y optimizar geométricamente modelos atómicos de interfaces Heusler/MXene, evaluando su estabilidad termodinámica y la energía de adhesión interfacial.
    \item Determinar la estructura de bandas electrónicas, la densidad de estados (DOS) y la polarización magnética local en la interfaz.
    \item Evaluar el efecto del acoplamiento espín-órbita (SOC) en la dispersión de bandas y calcular los invariantes topológicos (como el número de Chern o la Z2) para clasificar las fases electrónicas del sistema.
    \item Simular la influencia de estímulos externos (tensión mecánica epitaxial y campos eléctricos transversales) sobre la modulación de las propiedades topológicas y magnéticas de las heteroestructuras.
\end{enumerate}
